% Section 4 -- Lane-Matching Compiler

\subsection{Compiler Overview}\label{sec:compiler:overview}

The lane-matching compiler turns a single-GPU COVER/Hybrid-Cover
schedule into a multi-lane Stateflow plan. Its input is a list of
\emph{buffer states} -- each state is a $q$-subset of partitions that
would occupy one device for one microstep on a single GPU -- together
with the per-state edge buckets that cover the resulting edges. Its
output is a permutation of those states arranged into rounds of
$k$ lanes, along with a populated \texttt{StateflowPlan} including
cross-lane \texttt{PEER\_RELAY} descriptors.

The compiler runs in three passes (Figure~\ref{fig:compiler:pipeline}).
First, \emph{lane assignment} chooses, for each round, which buffer
states run on which lane so as to maximize residency overlap with the
previous round. Second, \emph{plan construction} materializes the
permutation into a \texttt{StateflowPlan} with per-lane microstates and
intra-lane handoffs. Third, \emph{cross-lane lifting} walks the plan
and emits a \texttt{PeerHandoffDescriptor} whenever a microstate
admits a partition that a peer lane held in the previous round.

\begin{figure}[t]
\centering
\begin{lstlisting}[basicstyle=\ttfamily\footnotesize]
buffer_states, edge_buckets  (single-GPU COVER output)
        |
        v
[Pass 1] assign_lanes_per_round  -> permutation: [s0,s1,s2,s3,...]
        |                            (round r uses permutation[r*k..r*k+k])
        v
[Pass 2] build_multi_gpu_stateflow_plan_from_permutation
        |   ->  plan with per-lane microstates + intra-lane handoffs
        v
[Pass 3] populate_cross_lane_handoffs
            ->  plan.cross_lane_handoffs += PEER_RELAY descriptors
        |
        v
[Pass 4] score_stateflow_plan    (selection_cost, planned_peer_bytes)
        |
        v
[Pass 5] validator (V1--V5)      (reject or accept)
\end{lstlisting}
\caption{Lane-matching compiler pipeline. Passes 1--3 produce the
plan; Pass 4 scores it; Pass 5 rejects it on invariant violation.}
\label{fig:compiler:pipeline}
\end{figure}

The rest of this section fills in the two non-trivial passes. Passes
2, 4, and 5 are mechanical: Pass~2 re-packs tensors into microstates
(Section~\ref{sec:ir:types}); Pass~4 is Equation~\ref{eq:cost}; Pass~5
is the validator of Section~\ref{sec:ir:validator}.

\subsection{Pass 1: Per-Round Lane Assignment}\label{sec:compiler:assign}

At each round $r$, the compiler must choose $k$ buffer states from
those not yet placed, and it must order them across lanes
$\ell \in \{0, \dots, k-1\}$. Two states may be placed in the same
round only if they are \emph{disjoint} -- the intersection of their
partition sets is empty -- otherwise they would double-load the same
partition on two devices and break superstate disjointness (invariant
V4). Compatibility is precomputed as a symmetric boolean matrix over
all state pairs.

Given disjointness as a hard constraint, the assignment problem at
round $r$ is:
\[
\arg\min_{(a_0, \dots, a_{k-1})}
  \underbrace{\sum_{\ell=0}^{k-1}
              T(\mathrm{prev}_\ell,\, a_\ell)}_{\text{transition cost}}
  \ +\
  \underbrace{I(a_0, \dots, a_{k-1})}_{\text{imbalance penalty}},
\]
where $a_\ell$ is the state assigned to lane $\ell$,
$\mathrm{prev}_\ell$ is what lane $\ell$ held in round $r\!-\!1$, and
the sum runs over all disjoint $k$-tuples drawn from the remaining
states.

\paragraph{Transition cost.} The per-lane term $T(p, n)$ prices the
handoff from previous state $p$ to next state $n$ on the same lane:
\begin{equation}
T(p, n) = \frac{\mathrm{host\_bytes}(p, n)}{B_{\mathrm{host}}}
         + [\mathrm{overlap}(p, n) = 0] \cdot w_{\mathrm{bnd}}.
\label{eq:transition}
\end{equation}
Here $\mathrm{host\_bytes}(p, n)$ is the byte volume that would need
to be loaded from host if the admit were served by \texttt{FULL\_RELOAD}
(partitions in $n$ but not in $p$, scaled by the embedding layout);
$B_{\mathrm{host}}$ is the host bandwidth from
\texttt{LaneMatchCostConfig}; and $w_{\mathrm{bnd}}$ is an additive
boundary penalty applied whenever residency changes completely between
rounds (the typical case for disjoint-round baselines).

Equation~\ref{eq:transition} is intentionally \emph{host-denominated}.
In Phase 3, peer transfers are planned but not executed, so crediting
a lane assignment with peer-bandwidth savings it will not realize
would make the compiler over-eager. A configuration flag
(\texttt{allow\_peer\_relay}) enables an optional $-\,\mathrm{peer\_bytes}
/B_{\mathrm{peer}}$ credit term; it is disabled by default and guarded
behind \texttt{GEGE\_STATEFLOW\_ALLOW\_PEER\_RELAY=1}. We leave it off
for the headline measurements.

\paragraph{Imbalance penalty.} The per-round term $I$ discourages
assignments where one lane is handed a large bucket and the other a
small one, which would waste the gap between their respective
completions:
\begin{equation}
I(a_0, \dots, a_{k-1}) = w_{\mathrm{imb}}
  \cdot \frac{\sum_\ell (E_{a_\ell} - \bar E)^2}{\bar E},
\label{eq:imbalance}
\end{equation}
where $E_{a_\ell}$ is lane $\ell$'s total edge-bucket volume for the
round and $\bar E$ is the mean across lanes. This is a chi-square-style
normalization: the penalty scales with relative, not absolute,
imbalance. $w_{\mathrm{imb}}$ defaults to $0.25$.

\subsection{Solvers}\label{sec:compiler:solvers}

The per-round problem has $\binom{m}{k} \cdot k!$ candidate
$k$-tuples, where $m$ is the number of unplaced states. Three solvers
are provided; the choice is controlled by
\texttt{GEGE\_STATEFLOW\_LANE\_MATCH\_SOLVER}.

\begin{description}
\item[\texttt{optimal2} (default).] Historical name for the exact
small-$k$ solver. The compiler enumerates all compatible disjoint
$k$-tuples and, for each tuple, solves the lane-assignment subproblem
exactly with subset DP over the $k \times k$ transition-cost matrix.
At $k{=}2$ this degenerates to the original exhaustive pair search;
at $k{=}4$ it remains easily tractable and preserves the same
host-bytes / boundary / imbalance objective as the 2-lane case.

\item[\texttt{greedy}.] Incremental search that places states one at a
time, choosing each to maximize overlap with the corresponding
previous-round lane. Cheap, but it can commit early to an assignment
that is locally good and globally suboptimal. Used as a fallback when
the exact solver is intentionally disabled.

\item[\texttt{hungarian} (alias for \texttt{optimal2}).] Accepted for
backward compatibility with earlier drafts. The current implementation
routes it to the same exact small-$k$ solver.
\end{description}

At two GPUs -- the focus of this paper -- exhaustive enumeration is
both optimal and trivially auditable. The same exact objective now
extends unchanged to four GPUs; the only difference is that the
assignment step is solved with subset DP instead of a hand-written
two-lane special case. \texttt{optimal2} remains the default.

\subsection{Pass 3: Cross-Lane Lifting}\label{sec:compiler:lift}

Once the per-round assignment is baked into a permutation and
materialized by Pass 2, the resulting plan already has correct
per-lane handoffs (each \texttt{LanePlan} has $N$ microstates and $N-1$
intra-lane \texttt{HandoffPlan}s). It does \emph{not} yet know which
admits could have been peer-served; that information is computed in
Pass 3 by walking the plan.

\begin{algorithm}[t]
\caption{Cross-lane lifting (\texttt{populate\_cross\_lane\_handoffs})}
\label{alg:lift}
\begin{algorithmic}[1]
\State $\text{cross\_lane\_handoffs} \gets \emptyset$
\For{each lane $\ell$ and each microstate $m$ on $\ell$ with $m \ne m_0$}
  \State let $m_{-1}$ be the previous microstate on $\ell$
  \For{each partition $\pi \in m.\mathrm{resident\_partitions}$ not in $m_{-1}$}
    \State $\text{best} \gets \bot$
  \For{each other lane $\ell'$}
      \State let $m'$ be $\ell'$'s microstate at superstate $m.\mathrm{superstate\_id}{-}1$
      \If{$m' \ne \bot$ and $\pi \in m'.\mathrm{resident\_partitions}$}
        \State $b \gets \mathrm{partition\_transfer\_bytes}(\pi)$
        \If{$\text{best} = \bot$ \textbf{or} $\ell' < \text{best.lane}$}
          \State $\text{best} \gets (\ell', m', b)$
        \EndIf
      \EndIf
    \EndFor
    \If{$\text{best} \ne \bot$}
      \State emit \texttt{PEER\_RELAY} handoff from $\text{best}.\ell'$ to $\ell$, partition $\pi$
      \State $\text{plan.total\_peer\_handoff\_bytes} \mathrel{+}= \text{best.bytes}$
    \EndIf
  \EndFor
\EndFor
\end{algorithmic}
\end{algorithm}

Algorithm~\ref{alg:lift} is complete: for every admit on every lane,
it either finds a peer source (in which case a \texttt{PeerHandoff}
descriptor is emitted and accounted against
\texttt{total\_peer\_handoff\_bytes}) or it does not (in which case
the admit will be served by host fallback at runtime, already represented
by the intra-lane \texttt{FULL\_RELOAD} in Pass 2). Deduplication by
$(\ell, m, \pi)$ prevents double-counting when a partition is
simultaneously resident on multiple peer lanes.

\paragraph{Source tie-breaking.} Under the current byte model,
\texttt{partition\_transfer\_bytes} depends only on the partition being
transferred, not on which peer lane serves it. When two peer lanes
could each service the same admit, we therefore break ties by smaller
lane id. The rule is deliberately simple and deterministic; a future
bandwidth-aware model could replace it with link-specific source
selection without changing the rest of the pass.

\paragraph{What Pass 3 does \emph{not} do.} It does not modify the
lane assignment chosen by Pass 1. The per-lane permutation is frozen
by the time Pass 3 runs. Co-designing the two passes -- letting the
availability of a peer source pull the per-lane assignment toward
configurations that maximize such sources -- is a natural next step.
The cleaner staging is to run them independently, confirm the per-lane
story is correct, then extend.

\subsection{Variants and Baseline}\label{sec:compiler:variants}

The compiler produces two named multi-GPU variants, both of which
share Passes 2--5:

\begin{description}
\item[\textsc{DISJOINT\_ROUNDS}.] A baseline that skips Pass 1 entirely
and routes to \texttt{getDisjointBufferStatePermutation}: states are
partitioned into disjoint rounds in their COVER-enumeration order,
then assigned to lanes in round order. This is the behavior of running
two independent COVER schedulers on alternating rounds -- the status
quo for multi-GPU COVER-based trainers (Section~\ref{sec:intro}).
\item[\textsc{LANE\_MATCHED}.] The full Pass 1 + Pass 3 pipeline
described above. This is the variant the compiler would choose under
default scoring on Twitter (Section~\ref{sec:eval}).
\end{description}

Both variants satisfy V1--V5; both populate
\texttt{cross\_lane\_handoffs} and \texttt{total\_peer\_handoff\_bytes}.
The difference is entirely in the lane assignment. On Twitter 16p
2-GPU, this difference alone changes the handoff-mode distribution
from $\{\texttt{FULL\_RELOAD}:10,\ \texttt{ROTATING\_OVERWRITE}:8,\
\texttt{PEER\_RELAY}:8\}$ to
$\{\texttt{FULL\_RELOAD}:0,\ \texttt{ROTATING\_OVERWRITE}:18,\
\texttt{PEER\_RELAY}:18\}$, eliminating every cold admit from the plan.

\subsection{Pass 4: Scoring}\label{sec:compiler:score}

Scored plans expose every term of Equation~\ref{eq:cost} on
\texttt{StateflowPlan.cost\_breakdown}:
\texttt{admitted\_partition\_cost},
\texttt{bucket\_edge\_cost},
\texttt{boundary\_cost},
\texttt{lane\_imbalance\_cost}, and the summed
\texttt{selection\_cost}. For multi-GPU plans, the admitted term is
byte-accurate (summed $\mathrm{rows}(\pi) \cdot \mathrm{bytes\_per\_row}$
over admitted objects); for single-GPU plans, it falls back to a
count-based weighted admission load (preserving backward compatibility
with existing single-GPU cost reports).

\texttt{planned\_peer\_bytes} is populated alongside
\texttt{selection\_cost} but \emph{not summed into it}. This is the
same split introduced in Section~\ref{sec:ir:cost}: the compiler
reports the Phase 4 opportunity without letting itself claim credit
for it. A consumer that wants to optimize under a peer-aware cost
model can sum the two terms explicitly; the default pipeline does not.

\subsection{Complexity and Determinism}\label{sec:compiler:complexity}

Per round, Pass 1 does
$O\!\left(\binom{m}{k} \cdot k \cdot 2^k\right)$ work with the exact
\texttt{optimal2} solver and reduces to $O(m^2)$ at $k{=}2$;
Pass 3 does $O(k \cdot N \cdot q)$ work to scan candidate peer sources
(where $N$ is the number of microstates per lane and $q$ is the buffer
capacity). Pass 4 is linear in the plan size. At Twitter 16p / 2-GPU,
the entire compile takes a few milliseconds and is completely
dominated by I/O overhead elsewhere in the training pipeline.

Every solver uses total orderings on ties (state-index, then lane-id),
so plans are deterministic for a fixed input. This matters for
reproducibility in the evaluation (Section~\ref{sec:eval}) and for the
validator: V1--V5 are strict equalities, so a non-deterministic
compiler would randomly fail otherwise-valid runs.
